% ==============================================================================
% INTRODUCTION GÉNÉRALE
% ==============================================================================
\chapter*{Introduction Générale}
\addcontentsline{toc}{chapter}{Introduction Générale}
\markboth{INTRODUCTION GÉNÉRALE}{}

\vspace{0.3cm}

\begin{infobox}[Contexte de l'étude]
L'essor du numérique dans l'enseignement supérieur génère chaque année
des millions de retours étudiants. Traiter ces données manuellement est
devenu impossible. L'intelligence artificielle offre une réponse concrète.
\end{infobox}

\vspace{0.5cm}

\section*{Contexte général}

L'enseignement supérieur moderne est caractérisé par un flux continu et
croissant d'interactions entre enseignants et étudiants. Parmi celles-ci,
les commentaires et retours d'expérience sur les cours constituent une
source d'information précieuse pour améliorer la qualité pédagogique des
enseignements. Cependant, la multiplication des effectifs étudiants et la
diversification des modalités d'enseignement — notamment en ligne — ont
rendu le traitement manuel de ces retours pratiquement impossible à grande
échelle.

Face à ce constat, l'analyse automatique de sentiments (\textit{Sentiment
Analysis}) émerge comme une solution technologique particulièrement
prometteuse. Discipline à l'intersection du \acrfull{taln}, de l'apprentissage
automatique et de la linguistique computationnelle, elle permet d'identifier
et de classifier automatiquement les opinions exprimées dans un texte, sans
intervention humaine.

\section*{Problématique}

La langue française présente des défis spécifiques pour l'analyse
automatique de sentiments : richesse morphologique, usage fréquent de la
négation, expressions idiomatiques, niveaux de registre variés et présence
d'accents. Ces particularités rendent inadéquates les approches développées
pour l'anglais, et nécessitent des modèles spécialement adaptés à la langue
française.

Comment peut-on concevoir un système robuste, précis et accessible permettant
d'analyser automatiquement et en temps réel les sentiments exprimés dans les
commentaires d'étudiants en français, et de mettre ces analyses à la
disposition des enseignants et des administrateurs de manière exploitable ?

\section*{Solution proposée}

Ce projet propose une réponse complète à cette problématique à travers
le développement d'un système multi-couches intégrant :

\begin{itemize}[leftmargin=2em]
  \item[\textcolor{bleuprincipal}{\ding{118}}]
    Un modèle d'\acrfull{ia} basé sur \textbf{CamemBERT}, architecture
    Transformer pré-entraînée sur des milliards de tokens français, adaptée
    par fine-tuning à la classification de sentiments étudiants en trois
    catégories.

  \item[\textcolor{vertprincipal}{\ding{118}}]
    Un \textbf{serveur API RESTful} développé avec FastAPI, assurant
    l'interface entre le modèle \acrshort{ia} et les applications
    clientes, avec persistance des données dans PostgreSQL.

  \item[\textcolor{orangeaccent}{\ding{118}}]
    Un \textbf{tableau de bord analytique} Streamlit permettant aux
    enseignants de suivre en temps réel les tendances de sentiment par
    cours et par matière.

  \item[\textcolor{rougeprincipal}{\ding{118}}]
    Une \textbf{application mobile Android} Flutter offrant aux étudiants
    une interface simple et moderne pour soumettre leurs commentaires depuis
    n'importe où.
\end{itemize}

\section*{Objectifs du projet}

L'objectif général de ce \acrfull{pfe} est de concevoir, développer et
valider un système complet d'analyse de sentiments en temps réel pour les
commentaires étudiants en langue française. De manière plus spécifique,
ce projet vise à :

\begin{enumerate}[leftmargin=2.5em]
  \item Constituer et annoter un jeu de données de commentaires étudiants
        en français, équilibré entre les classes de sentiment.
  \item Fine-tuner le modèle CamemBERT pour la classification ternaire de
        sentiments avec des performances acceptables.
  \item Développer une API REST robuste et documentée exposant les
        fonctionnalités du modèle.
  \item Concevoir un tableau de bord analytique interactif pour la
        visualisation des résultats.
  \item Développer une application mobile multiplateforme permettant la
        soumission et la consultation des analyses.
  \item Évaluer les performances du système et identifier les pistes
        d'amélioration.
\end{enumerate}

\section*{Organisation du mémoire}

Ce mémoire est organisé en cinq chapitres, précédés d'une introduction et
suivis d'une conclusion générale :

\vspace{0.3em}
\noindent\textbf{Chapitre 1 — Contexte et Problématique :} présente le
contexte général de l'enseignement supérieur numérique, analyse les limites
des approches existantes et justifie les choix technologiques retenus.

\vspace{0.3em}
\noindent\textbf{Chapitre 2 — Analyse et Conception :} détaille l'analyse
fonctionnelle et non fonctionnelle du système, présente les diagrammes
\acrshort{uml} et l'architecture globale, et justifie les choix techniques.

\vspace{0.3em}
\noindent\textbf{Chapitre 3 — Implémentation :} décrit en détail les
environnements de développement, la structure du projet, les algorithmes
clés et les extraits de code commentés des composants principaux.

\vspace{0.3em}
\noindent\textbf{Chapitre 4 — Tests et Résultats :} présente la stratégie
de validation, les métriques de performance du modèle, les scénarios de
test fonctionnels et l'analyse critique des résultats obtenus.

\vspace{0.3em}
\noindent\textbf{Chapitre 5 — Discussion et Perspectives :} propose une
analyse critique du projet, identifie ses limites actuelles et trace les
perspectives d'évolution futures.

\cleardoublepage
